\documentclass[a4paper,10pt]{article}
\usepackage[left=3cm,top=3cm,right=2.5cm,bottom=4cm]{geometry} 
\usepackage[utf8]{inputenc}
\usepackage[activeacute,spanish]{babel}
\usepackage{multirow}	% Para poder unir filas en las tablas
\usepackage{colortbl}	% Para colorear tablas
\usepackage{latexsym,amssymb,oldlfont,amsmath}
\usepackage{graphicx,graphics}
\usepackage{multicol}
\usepackage{color}
\usepackage{listings}
\usepackage{lastpage}
\usepackage{enumitem}
\usepackage{dsfont} %para poner numero reales


%\usepackage{fontspec}
%\setmainfont{Times New Roman}
\newcommand{\R}{{\mathbb R}}
\newcommand{\si}{{\mathcal S}^n}
\renewcommand{\theequation}{\thesection.\arabic{equation}} % Cambia formato salida numeracion equation
\makeatletter\@addtoreset {equation}{section} \makeatother % Tambien se usa
\renewcommand{\baselinestretch}{1.3}
\usepackage{amssymb}
\usepackage{hyperref}
\definecolor{midnightblue}{rgb}{0.1,0.1,0.44}


\newcommand{\dis}{\displaystyle}
\renewcommand{\theenumii}{\arabic{enumii}}

\def\Z{\mathbb{Z}}
\def\N{\mathbb{N}}
\def\R{\mathbb{R}}
\def\M{\mathcal{M}_n(\mathbb{R})}
\def\sin{\operatorname{sen}}
\def\e{\operatorname{e}}
\DeclareMathOperator{\arcsec}{arcsec}
\DeclareMathOperator{\arccot}{arccot}
\DeclareMathOperator{\arccsc}{arccsc}



\usepackage{etoolbox,fancyhdr,xcolor}
\fancyhf{} % sets both header and footer to nothing
\renewcommand{\headrulewidth}{0pt}
%\renewcommand{\footrulewidth}{0.5pt}
\pagestyle{fancy}
\headheight=60pt %para cambiar el tamaño del encabezadox
\fancyhead[L] %la "L" indica a la izquierda
{
\begin{minipage}{6cm}
\includegraphics[width=0.9\textwidth]{logo_udp.eps}
\end{minipage}
}

%\fancyhead[C]
%{
%\textbf{Cálculo II - Guía 1}
%}


%\fancyhead[R] %la "L" indica a la izquierda
%{

%}
%\begin{minipage}{4cm}
%\small %tamaño de letra pequeño
%{
%\textsf
%{
%Universidad de Chile\\
%...}}
%\end{minipage}}



%\fancyfoot[C]{}
%{
%\begin{minipage}{10cm}
%\color{midnightblue}
%\begin{center}Casa Matriz: Monjitas 527 Piso 15 - Santiago de Chile\\
%Fono (02)24367255 \end{center}
%\end{minipage}
%}
\fancyfoot[C]{}
%Página \thepage\ de \pageref{LastPage}}% Páginas Pares numeradas a la derecha
%\fancyfoot[LO]{\thepage} % Páginas Impares numeradas a la izquierda

\title{C2_Guía de Ejercicios de verdaderos o falsos}

\begin{document}

\begin{center}
Cálculo II\\
Guía de Ejercicios de verdaderos o falsos\\
\end{center}
\vspace{10mm}

\begin{enumerate}


\item Determine si las siguientes afirmaciones son verdaderas o falsas. Justifique.
\begin{enumerate}[label=\emph{\alph*)}]
\item $\dis \int \cos(kx)dx = \frac{1}{k}\sin(kx)+C$ con $k \in R - \left\{ 0\right\}$, $C \in R$
\item $\dis \int \sin(kx)dx = -\frac{1}{k}\cos(kx)+C$ con $k \in R - \left\{ 0\right\}$, $C \in R$
\end{enumerate}

\item Determine si las siguientes afirmaciones son verdaderas o falsas. Justifique.
\begin{enumerate}[label=\emph{\alph*)}]
\item $\dis \int f(x) g(x)\, dx = \dis \int f(x)dx \int g(x)dx$
\item $\dis \int f(x)dx(\int g(x)dx + \int h(x)dx) = \dis  \int f(x)dx\int g(x)dx + \int f(x)dx\int h(x)dx$
\item $\dis \int (f(x) +  g(y))dx = \dis  \int f(x)dx +  g(y)x$
\item $\dis \int g'(z) f(g(z))\, dz = \dis \int f(t)dt$
\item $\dis \int xh(t)g(x^2)\, dx = \frac{h(t)}{2}\dis \int f(z)dz$
\item $\dis \int 3x^2 \ln(x^3)f(t)\, dt = f(t)\dis \int \ln(u)du$
\end{enumerate}

\item  Determine si las siguientes afirmaciones son verdaderas o falsas. Justifique.
\begin{enumerate}[label=\emph{\alph*)}]
\item Un fabricante descubrió que el costo marginal cuando se producen $q$ unidades es \\ $3q2 -60q+400$ dólares por unidad. Si el costo total de producción de las 2 primeras unidades es $U\$900$, entonces  el costo total de producción de las 3 primeras unidades es $U\$1350$.
\item Una compañía actualmente produce $x$ unidades por semana de producto. Por experiencia, saben que el costo de producir una unidad más cuando la producción es $x$ unidades está dada por la función costo marginal $C'(x)= 25-0.02x$, entonces de acuerdo a ese costo marginal  el costo extra para elevar la producción de 150 a 200 unidades  es $\$1075$. 
\item El ingreso marginal de una empresa está dado por $R'(x) =15-0,01x$ donde $x$ es cantidad de artículos vendidos. Suponiendo que el ingreso es cero cuando no se venden unidades y que cada producto que produce y vende la empresa  se vende a un precio $p$, entonces la relación entre el precio $p$ y la cantidad $x$ es $p= 15-0,005x$.

\end{enumerate}

\end{enumerate}
\vspace{15mm}

Edición de la guía: Sara Arancibia y Rodrigo Ábel
\end{document}